\chapter{Introduction}

The production of 2D animation still depends heavily on manual, frame-by-frame motion refinement. This makes dense correspondence estimation a particularly valuable primitive for modern anime pipelines, where optical flow can support frame interpolation, inbetweening, color propagation, temporally consistent editing, and downstream video synthesis~\cite{animeinterp,superSloMo,softmaxSplat,linedrawingFlow,toonTracking}. In practice, however, animation is a hostile domain for optical flow: large flat-color regions weaken local matching, contour lines create sharp motion discontinuities, occlusions are frequent, and artist-driven deformations often violate the assumptions that work well in natural videos.

These properties create a serious gap between natural-scene optical flow and animation-oriented correspondence estimation. While supervised methods such as PWC-Net, RAFT, GMFlow, and FlowFormer have driven major progress on benchmarks like Sintel and FlyingThings3D~\cite{pwcnet,raft,gmflow,flowformer,sintel,ft3d}, their design biases are not automatically well aligned with cel-style animation. The AnimeRun benchmark~\cite{animerun} makes this mismatch explicit: strong natural-image models still struggle on texture-sparse interiors, layered occlusion boundaries, and very large stylized motion.

The lack of dense real-animation labels further increases the difficulty. Large-scale supervision for stylized content is expensive to build, hard to diversify, and not representative of the full range of production animation. This motivates unsupervised optical flow learning, but standard unsupervised objectives are also brittle in animation because brightness constancy is often weak, flat regions are ambiguous, and pairwise matching alone is rarely sufficient to stabilize object identity over time.

This thesis addresses that gap through \textbf{AniUnFlow}, an unsupervised object-memory optical-flow framework for 2D animation. AniUnFlow organizes motion estimation around SAM-derived object slots, temporal slot memory, dense refinement, and a propagation-memory signal obtained by warping previous SAM structure into the current frame. Rather than treating segmentation, dense matching, and temporal context as separate cues, AniUnFlow uses them as cooperating priors: object structure stabilizes the motion hypothesis, dense matching restores local detail, and propagated SAM agreement constrains refinement where correspondence is trustworthy.

The thesis also documents the design evolution that led to the current framework. Earlier internal branches established the value of structure-first motion reasoning, while the latest large-motion extension broadens the same line of work toward non-rigid slot flow, coarse global search, and visibility-aware compositing. In the body of this thesis, however, the writing remains centered on a single unified framework name: \textbf{AniUnFlow}.

\textbf{The main contributions of this thesis are as follows:}
\begin{itemize}
\item We present \textbf{AniUnFlow}, an unsupervised optical-flow framework for 2D animation built around SAM-guided object memory, dense refinement, boundary-aware correction, and propagated structural memory.
\item We introduce a \textbf{SAM propagation memory} mechanism that warps neighboring SAM structure forward in time and converts its agreement with the current frame into a dense reliability signal for confidence modulation, refinement, and regularization.
\item We show that the same structure-first design also enables a practical \textbf{interpretability perspective} for animation flow, because the framework exposes readable intermediate signals such as slot priors, propagated structural agreement, confidence maps, and residual correction activity.
\item We reorganize the optical-flow narrative for animation around a clear \emph{structure-first, dense-refine} principle and provide richer paper-oriented visualization, ablation-style decomposition, and branch-level analysis, including a concise discussion of the latest large-motion extension.
\end{itemize}
